% Phần lý thuyết của báo cáo. Ghép sau abstract và keywords.
% Yêu cầu các gói toán thông dụng amsmath và amssymb.
% Không chứa documentclass, preamble, tài liệu tham khảo hoặc số liệu đo.

\section{Giới thiệu}
\label{sec:theory-introduction}

Bài toán $N$ quân hậu (\emph{N-Queens}) yêu cầu đặt $N$ quân hậu trên bàn cờ
$N\times N$ sao cho không có hai quân tấn công nhau. Cách phát biểu ngắn gọn
nhưng nhiều nhóm ràng buộc giao nhau khiến bài toán phù hợp để khảo sát ảnh
hưởng của cách biểu diễn đối với bộ giải ràng buộc. Trong cách tiếp cận SAT,
mỗi ô bàn cờ được biểu diễn bằng một biến Boolean, còn các điều kiện về hàng,
cột và đường chéo được chuyển thành công thức ở dạng chuẩn hội
(\emph{conjunctive normal form}, CNF). Bộ giải SAT tìm một phép gán thỏa mãn;
phép gán này sau đó được giải mã và kiểm tra như một bàn cờ thực sự.

Báo cáo kế thừa hướng khảo sát mã hóa AMO trong \cite{bach2024,module3}
và triển khai năm phương án: pairwise, binary, sequential counter,
commander và product. Chúng biểu diễn cùng điều kiện ``nhiều nhất một'',
với những lựa chọn khác nhau về số biến phụ, số clause và cấu trúc suy
diễn. Hai yếu tố cần đánh giá đồng thời là chi phí tạo công thức và chi
phí tìm nghiệm trên công thức đó. PySAT cung cấp giao diện gọi SAT
solver \cite{pysat}; các mã hóa được xây dựng tường minh để đối chiếu
công thức, xuất CNF và kiểm tra kích thước mô hình.

Nghiên cứu tập trung vào ba câu hỏi: mỗi mã hóa làm thay đổi kích thước
CNF ra sao; thay đổi đó liên hệ như thế nào với thời gian giải khi tăng
kích thước bàn cờ; và hiệu năng SAT so với mô hình CP và MIP tương ứng
thay đổi như thế nào. Để trả lời, báo cáo sử dụng một mô hình N-Queens
thống nhất, chứng minh tính đúng của các mã hóa và kiểm tra độc lập
nghiệm trả về. Các mô hình đối chứng được xây dựng cho CP-SAT, CPLEX
CP, CPLEX MIP và Gurobi. Cách tổ chức này gắn phân tích lý thuyết với
một quy trình thực nghiệm có thể tái lập.

\section{Cơ sở bài toán và mã hóa SAT}
\label{sec:theory-background}

\subsection{Phạm vi bài toán N-Queens}

Xét bàn cờ trống với $N\geq 1$, các chỉ số hàng và cột thuộc
$\{0,\ldots,N-1\}$. Một nghiệm được lưu dưới dạng
$q=(q_0,\ldots,q_{N-1})$, trong đó $q_r$ là cột của hậu ở hàng $r$.
Các điều kiện hợp lệ là
\begin{align}
0\leq q_r<N &\qquad &&0\leq r<N, \label{eq:th-domain}\\
q_r\neq q_t &\qquad &&0\leq r<t<N, \label{eq:th-columns}\\
|q_r-q_t|\neq |r-t| &\qquad &&0\leq r<t<N. \label{eq:th-diagonal}
\end{align}
Biểu diễn này đã đặt đúng một hậu trên mỗi hàng. Điều kiện
\eqref{eq:th-columns} loại xung đột theo cột; điều kiện
\eqref{eq:th-diagonal} loại xung đột theo hai hướng đường chéo.

Phạm vi thực nghiệm là tìm \emph{một nghiệm} trên bàn trống hoặc kết
luận không tồn tại nghiệm. Biến thể \emph{N-Queens Completion} có thêm
các hậu cố định trong đầu vào, còn bài toán đếm yêu cầu xác định số
lượng tất cả nghiệm. Các nhiệm vụ này có phạm vi đánh giá khác nhau.
Trong báo cáo, tính đầy đủ NP (\emph{NP-completeness}) của SAT tổng
quát không được dùng làm phát biểu độ phức tạp cho N-Queens trên bàn trống.

\subsection{SAT và ràng buộc lực lượng}

Một biến Boolean nhận giá trị Đúng hoặc Sai. Một \emph{literal} là biến
hoặc phủ định của biến; một \emph{clause}, gọi là mệnh đề tuyển trong báo
cáo này, là phép OR của các literal. Công thức CNF là phép AND của các
clause. Chẳng hạn,
\begin{equation}
 (\neg x_1\lor x_2)\land(x_1\lor x_3)\land(\neg x_2\lor\neg x_3)
\end{equation}
là một công thức CNF. SAT (\emph{Boolean satisfiability}) hỏi có tồn tại
phép gán làm công thức đúng hay không. Kết quả SAT đi kèm một mô hình
thỏa mãn; UNSAT là kết luận không có phép gán thỏa mãn. Hết thời gian
không đồng nghĩa với UNSAT.

Với nhóm biến $X=(x_1,\ldots,x_m)$, ba ràng buộc lực lượng
(\emph{cardinality constraints}) được sử dụng là
\begin{align}
\operatorname{ALO}(X)&:\quad \sum_{i=1}^{m}x_i\geq 1,
&\operatorname{ALO}(X)&=\bigvee_{i=1}^{m}x_i,\label{eq:th-alo}\\
\operatorname{AMO}(X)&:\quad \sum_{i=1}^{m}x_i\leq 1,
&\operatorname{EXO}(X)&=\operatorname{ALO}(X)\land\operatorname{AMO}(X).
\label{eq:th-exo}
\end{align}
ALO (\emph{at least one}) nghĩa là ít nhất một biến đúng; AMO
(\emph{at most one}) nghĩa là nhiều nhất một biến đúng; EXO
(\emph{exactly one}) nghĩa là đúng một biến đúng. Trong tổng, giá trị Đúng
được hiểu là 1 và Sai là 0. ALO được biểu diễn bằng một clause trực tiếp.
Phần cần lựa chọn phương pháp mã hóa chủ yếu là AMO \cite{module3}.

Một mã hóa AMO có thể thêm tập biến phụ $Y$. Tiêu chuẩn tính đúng là
\begin{equation}
 \operatorname{AMO}(X)\quad\Longleftrightarrow\quad
 \exists Y\; E(X,Y),
 \label{eq:th-projection}
\end{equation}
nghĩa là mỗi phép gán hợp lệ trên biến gốc có ít nhất một cách mở rộng sang
biến phụ, và không phép gán không hợp lệ nào có cách mở rộng như vậy.
Đây là tính đồng khả thỏa dưới phép chiếu lên biến gốc
(\emph{equisatisfiability under projection}). Biến phụ không nhất thiết
có một giá trị duy nhất. Điều này đặc biệt quan trọng khi diễn giải các
biến trong binary và sequential counter.

\section{Phương pháp SAT cho N-Queens}
\label{sec:theory-method}

\subsection{Biến và ràng buộc bàn cờ}

Đặt $x_{r,c}$ đúng khi và chỉ khi có một hậu tại ô $(r,c)$. Biến ô được
gán số nguyên dương
\begin{equation}
 \operatorname{id}(r,c)=rN+c+1.
 \label{eq:th-id}
\end{equation}
Như vậy $N^2$ biến đầu là biến bàn cờ, còn biến phụ được cấp phát liên tục
từ $N^2+1$. Mỗi lời gọi bộ mã hóa AMO nhận tập biến phụ riêng; việc dùng
chung bộ cấp số bảo đảm không trùng biến giữa các hàng, cột và đường chéo.

Các nhóm biến được xác định bởi
\begin{align}
 R_r &= \{x_{r,c}:0\leq c<N\},
 &C_c &= \{x_{r,c}:0\leq r<N\},\\
 D^-_d &= \{x_{r,c}:r-c=d\},
 &D^+_s &= \{x_{r,c}:r+c=s\}.
\end{align}
Trong hai tập đường chéo, các ô luôn thỏa $0\leq r,c<N$;
$1-N\leq d\leq N-1$ và $0\leq s\leq 2N-2$.
Công thức bàn cờ trước khi thay AMO bằng CNF là
\begin{equation}
\begin{split}
 F_N={}&\bigwedge_{r=0}^{N-1}\operatorname{EXO}(R_r)
 \land\bigwedge_{c=0}^{N-1}\operatorname{EXO}(C_c)\\
 &\land\bigwedge_{d=1-N}^{N-1}\operatorname{AMO}(D^-_d)
 \land\bigwedge_{s=0}^{2N-2}\operatorname{AMO}(D^+_s).
\end{split}
\label{eq:th-board}
\end{equation}
Chỉ hàng và cột có ALO; các đường chéo chỉ có AMO. AMO trên nhóm có
không quá một phần tử luôn đúng nên không sinh clause. Mặc dù ALO cột có
thể suy ra từ EXO hàng và AMO cột, cài đặt vẫn giữ nó trong cả năm
phương án để các phương án dùng cùng mô hình cấu trúc. Không phương án
nào được bổ sung ràng buộc phá đối xứng viết tay.

\subsection{Mã hóa pairwise}

Pairwise, còn gọi là binomial, cấm mọi cặp biến cùng đúng:
\begin{equation}
 E_{\mathrm{pw}}(X)=
 \bigwedge_{1\leq i<j\leq m}(\neg x_i\lor\neg x_j).
 \label{eq:th-pairwise}
\end{equation}
Phương án này không cần biến phụ và tạo $\binom{m}{2}$ clause.
Nếu có hai biến đúng thì clause của cặp đó sai. Ngược lại, nếu có
nhiều nhất một biến đúng, mọi cặp đều chứa ít nhất một biến sai nên mọi
clause đều đúng. Vì vậy \eqref{eq:th-pairwise} biểu diễn chính xác AMO.
Ưu điểm là công thức đơn giản; hạn chế là số clause tăng bậc hai theo
kích thước một nhóm.

\subsection{Mã hóa binary}

Đặt $k=\lceil\log_2m\rceil$ và thêm $k$ biến phụ $b_0,\ldots,b_{k-1}$.
Biến $x_i$ được gắn mã nhị phân của $i-1$, với bit thấp nhất được xét
trước, đúng theo thứ tự trong chương trình. Đặt $\ell_{i,j}=b_j$ nếu
bit $j$ của mã bằng 1, và $\ell_{i,j}=\neg b_j$ nếu bit bằng 0. Mã hóa là
\begin{equation}
 E_{\mathrm{bin}}(X,b)=
 \bigwedge_{i=1}^{m}\bigwedge_{j=0}^{k-1}(\neg x_i\lor\ell_{i,j}).
 \label{eq:th-binary}
\end{equation}
Mỗi biến đúng ép các bit phụ nhận đúng mã của nó. Hai biến phân biệt có
mã khác nhau tại ít nhất một bit, nên không thể cùng đúng. Nếu chỉ có
$x_i$ đúng, chọn các bit theo mã của $i-1$; nếu tất cả biến gốc sai, chọn
các bit bất kỳ. Hai chiều của \eqref{eq:th-projection} vì thế được bảo đảm.

Binary tạo $k$ biến phụ và $mk$ clause. Các liên kết chỉ có chiều
$x_i\rightarrow\operatorname{code}(i-1)$; không ép chiều ngược.
ALO, khi cần, được thêm trực tiếp trên các biến gốc. Cách này vẫn đúng
khi $m$ không phải lũy thừa của 2 và một số mã bit không đại diện cho
biến gốc nào.

\subsection{Mã hóa sequential counter}

Sequential counter, tức bộ đếm tuần tự, dùng $m-1$ biến phụ
$s_1,\ldots,s_{m-1}$ theo dạng AMO của Sinz \cite{sinz}. Cài đặt sử dụng
đúng các clause sau với mọi $m\geq2$:
\begin{align}
 &\neg x_1\lor s_1,\label{eq:th-seq-first}\\
 &\neg x_i\lor s_i &&2\leq i\leq m-1,\label{eq:th-seq-set}\\
 &\neg s_{i-1}\lor s_i &&2\leq i\leq m-1,\label{eq:th-seq-prop}\\
 &\neg x_i\lor\neg s_{i-1} &&2\leq i\leq m-1,\label{eq:th-seq-stop}\\
 &\neg x_m\lor\neg s_{m-1}.\label{eq:th-seq-last}
\end{align}
Các biểu thức trong miền $2\leq i\leq m-1$ được bỏ qua khi $m=2$.
Do đó số biến phụ luôn là $m-1$ và số clause là $3m-4$;
chương trình không thay các nhóm nhỏ bằng pairwise trong phương án này.

Để chứng minh chiều loại trừ nghiệm sai, giả sử $x_i$ và $x_j$ cùng
đúng với $i<j$. Biến $x_i$ làm $s_i$ đúng; các clause truyền trạng thái
làm $s_{j-1}$ đúng. Clause \eqref{eq:th-seq-stop} hoặc
\eqref{eq:th-seq-last} khi đó cấm $x_j$, gây mâu thuẫn. Với chiều còn
lại, khi nhóm có nhiều nhất một biến đúng, đặt
$s_i=\bigvee_{t=1}^{i}x_t$. Phép gán này thỏa mọi clause trên.
Đó là một cách mở rộng hợp lệ; công thức không yêu cầu mọi mô hình
phải buộc $s_i$ tương đương hai chiều với OR của tiền tố.

\subsection{Mã hóa commander đệ quy}

Commander sử dụng biến đại diện nhóm để đưa AMO trên nhiều biến về AMO
trên các nhóm nhỏ. Với $m\leq3$, chương trình dùng pairwise. Với $m>3$,
chia dãy biến theo thứ tự thành $g=\lceil m/3\rceil$ nhóm $G_1,\ldots,G_g$;
mỗi nhóm có tối đa ba phần tử. Thêm commander $c_i$ cho nhóm $G_i$ và
các ràng buộc
\begin{align}
 &\operatorname{AMO}_{\mathrm{pw}}(G_i) &&1\leq i\leq g,\\
 &\neg x\lor c_i &&x\in G_i,\label{eq:th-comm-up}\\
 &\neg c_i\lor\bigvee_{x\in G_i}x &&1\leq i\leq g.
 \label{eq:th-comm-down}
\end{align}
Cuối cùng, gọi đệ quy cùng mã hóa commander cho $(c_1,\ldots,c_g)$.
Hai nhóm clause \eqref{eq:th-comm-up}--\eqref{eq:th-comm-down} biểu
diễn $c_i\leftrightarrow\bigvee_{x\in G_i}x$, kể cả nhóm cuối chỉ có
một hoặc hai biến. Không thêm ALO trên các commander.

Nếu hai biến cùng nhóm đúng, pairwise nội bộ loại chúng. Nếu chúng ở
hai nhóm khác nhau, hai commander tương ứng cùng đúng và bị AMO ở tầng
trên loại. Lập luận này quy nạp đến trường hợp cơ sở. Ngược lại, nếu nhóm
gốc có nhiều nhất một biến đúng, gán mỗi commander bằng OR của nhóm nó;
tầng tiếp theo cũng có nhiều nhất một commander đúng. Lặp lại phép gán
này tới gốc sẽ thỏa toàn bộ clauses, chứng minh tính đầy đủ.

Gọi $a_m$ và $c_m$ lần lượt là số biến phụ và clause của một AMO
commander. Với $m\leq3$, $a_m=0$ và $c_m=\binom{m}{2}$.
Với $m>3$, đặt $t_i=|G_i|$ thì
\begin{align}
 a_m&=g+a_g,\\
 c_m&=\sum_{i=1}^{g}\binom{t_i}{2}+m+g+c_g.
 \label{eq:th-comm-count}
\end{align}
Vì kích thước nhóm được giữ ở mức ba và số commander giảm xấp xỉ ba
lần qua mỗi tầng, tổng chi phí các tầng thuộc $O(m)$. Ngay tầng đầu
đã có $\lceil m/3\rceil$ biến phụ và $m$ clause liên kết, nên cả hai
đại lượng thực tế thuộc $\Theta(m)$. Phân tích này áp dụng cho biến thể
đệ quy đang cài đặt; biến thể một tầng với nhóm cỡ căn bậc hai có truy
hồi và kích thước khác.

\subsection{Mã hóa product hai chiều}

Với $m\leq3$, product trong chương trình dùng pairwise để tránh chi phí
biến phụ ở nhóm nhỏ. Với $m>3$, đặt
\begin{equation}
 p=\lceil\sqrt m\rceil,\qquad q=\left\lceil\frac{m}{p}\right\rceil.
\end{equation}
Gắn $x_i$ vào ô $\bigl(u(i),v(i)\bigr)$ của lưới $p\times q$, trong đó
\begin{equation}
 u(i)=\left\lfloor\frac{i-1}{q}\right\rfloor+1,
 \qquad v(i)=((i-1)\bmod q)+1.
\end{equation}
Thêm biến chọn hàng $r_1,\ldots,r_p$ và biến chọn cột $d_1,\ldots,d_q$.
Mỗi biến gốc liên kết tới đúng hai biến chọn:
\begin{equation}
 (\neg x_i\lor r_{u(i)})\land(\neg x_i\lor d_{v(i)}),
 \qquad 1\leq i\leq m.
 \label{eq:th-product-links}
\end{equation}
Hai tập biến chọn được ràng buộc bởi AMO pairwise. Nếu hai biến gốc khác
nhau cùng đúng, vị trí của chúng khác hàng hoặc khác cột trong lưới;
khi đó AMO tương ứng bị vi phạm. Nếu có nhiều nhất một biến gốc đúng,
có thể chọn chỉ hàng và cột của biến đó; nếu không có biến nào đúng,
chọn tất cả biến phụ sai. Do đó mã hóa thỏa \eqref{eq:th-projection}.

Với $m>3$, số biến phụ và clause là
\begin{equation}
 a_m=p+q,\qquad c_m=2m+\binom p2+\binom q2.
 \label{eq:th-product-count}
\end{equation}
Do lưới được chọn cân bằng, $a_m=\Theta(\sqrt m)$ và $c_m=\Theta(m)$.
Khi $pq>m$, một số ô không gắn biến gốc; chúng không gây vấn đề vì
chương trình chỉ thêm ALO trên biến gốc khi cần EXO. Không thay ALO này
bằng ALO của các biến chọn, vốn có thể cho phép chọn một ô trống.

\subsection{Tính đúng của mô hình và quy trình giải}

\noindent\textbf{Mệnh đề 1.}
Với mỗi $N\geq1$ và mỗi mã hóa AMO nêu trên, công thức CNF sinh từ
\eqref{eq:th-board} khả thỏa khi và chỉ khi tồn tại một lời giải N-Queens.

\noindent\emph{Chứng minh.}
Theo các lập luận ở từng mã hóa, mọi mô hình của CNF đều thỏa AMO trên
tập biến gốc của từng nhóm. ALO và AMO hàng buộc đúng một hậu trên mỗi
hàng, tức có tổng cộng $N$ hậu. Ràng buộc cột và hai họ đường chéo
cấm mọi cặp hậu tấn công nhau, nên mô hình giải mã là nghiệm hợp lệ.
Đây là tính đúng đắn theo chiều từ mã hóa sang bài toán (\emph{soundness}).
Ngược lại, lấy một bàn cờ hợp lệ và đặt các biến ô có hậu là đúng.
Các nhóm hàng, cột, đường chéo thỏa điều kiện tương ứng. Mỗi nhóm có
ít nhất một phép mở rộng sang biến phụ theo
\eqref{eq:th-projection}. Các tập biến phụ không giao nhau nên có thể
ghép các phép mở rộng thành một mô hình của toàn bộ CNF. Đây là tính
đầy đủ (\emph{completeness}).\hfill$\square$

Quy trình cài đặt được tóm tắt như sau:
\begin{enumerate}
 \item Kiểm tra $N$ là số nguyên dương và chọn một trong năm mã hóa.
 \item Cấp $N^2$ biến ô; duyệt các hàng và cột, thêm một ALO và một
       AMO cho mỗi nhóm; sau đó thêm AMO cho hai họ đường chéo.
 \item Nạp CNF vào SAT solver thông qua PySAT \cite{pysat} và gọi giải.
 \item Nếu SAT, chỉ lấy các biến dương có số từ 1 tới $N^2$, giải mã
       thành véc-tơ cột $q$ và kiểm tra bàn cờ độc lập với solver.
 \item Lưu trạng thái, nghiệm nếu có, số biến, số clause và các giai
       đoạn thời gian; phân biệt UNSAT với hết thời gian hoặc lỗi công cụ.
\end{enumerate}
Kiểm tra độc lập yêu cầu $q$ có đúng $N$ phần tử trong miền hợp lệ và
ba tập $\{q_r\}$, $\{r-q_r\}$, $\{r+q_r\}$ đều có lực lượng $N$.
Nếu solver trả SAT nhưng phép kiểm tra không đạt, chương trình báo lỗi
thay vì ghi nhận thành công. CNF cũng có thể xuất theo định dạng DIMACS,
trong đó mỗi clause kết thúc bằng 0, để dùng với một bộ giải bên ngoài.

\subsection{Kích thước công thức trên toàn bàn cờ}

Với $N\geq2$, sau khi bỏ các nhóm đường chéo một phần tử, có $6N-6$
nhóm AMO và tổng độ dài các nhóm là $4N^2-4$. Có $2N$ clause ALO.
Mỗi họ đường chéo có các độ dài $1,2,\ldots,N-1,N,N-1,\ldots,2,1$.
Vì vậy, nếu $a_m,c_m$ là số biến phụ và clause cho một AMO kích thước
$m$, tổng số biến và clause của cài đặt là
\begin{align}
 V(N)&=N^2+(2N+2)a_N+4\sum_{m=2}^{N-1}a_m,\label{eq:th-general-v}\\
 C(N)&=2N+(2N+2)c_N+4\sum_{m=2}^{N-1}c_m.\label{eq:th-general-c}
\end{align}
Các công thức này tính clauses trước khi bộ giải tiền xử lý; không khẳng
định rằng số clauses còn hoạt động bên trong solver giữ nguyên.

Đối với pairwise, thay $a_m=0$ và $c_m=\binom m2$ thu được
\begin{equation}
 V_{\mathrm{pw}}(N)=N^2,\qquad
 C_{\mathrm{pw}}(N)=\frac{N(N-1)(5N-1)}{3}+2N.
 \label{eq:th-pw-board-count}
\end{equation}
Đối với sequential counter đang cài đặt, không rút gọn nhóm kích thước
hai hoặc ba, thay $a_m=m-1$ và $c_m=3m-4$ thu được
\begin{align}
 V_{\mathrm{seq}}(N)&=5N^2-6N+2,\\
 C_{\mathrm{seq}}(N)&=12(N-1)^2+2N.
 \label{eq:th-seq-board-count}
\end{align}
Ví dụ, ở $N=4$, pairwise tạo 16 biến và 84 clauses; sequential tạo
58 biến và 116 clauses. Đây là phép đếm từ định nghĩa, không phải số
đo thời gian. Nó cho thấy mã hóa tăng trưởng tuyến tính trên mỗi nhóm
không nhất thiết có ít clauses hơn ở nhóm nhỏ.

Binary dùng $a_m=\lceil\log_2m\rceil$ và
$c_m=m\lceil\log_2m\rceil$ trong
\eqref{eq:th-general-v}--\eqref{eq:th-general-c}. Commander dùng truy
hồi \eqref{eq:th-comm-count}. Product dùng
\eqref{eq:th-product-count} khi $m>3$ và pairwise khi $m\leq3$.
Bảng~\ref{tab:th-growth} tổng hợp bậc tăng trưởng; số biến gốc $N^2$
được tách khỏi số biến phụ để tránh nhầm lẫn.

\begin{table}[htbp]
 \centering
 \caption{Bậc tăng trưởng của CNF theo đúng các biến thể đã cài đặt}
 \label{tab:th-growth}
 \begin{tabular}{lll}
  \hline
  Mã hóa & Biến phụ toàn bàn & Clauses toàn bàn\\
  \hline
  Pairwise & $0$ & $\Theta(N^3)$\\
  Binary & $\Theta(N\log N)$ & $\Theta(N^2\log N)$\\
  Sequential & $\Theta(N^2)$ & $\Theta(N^2)$\\
  Commander đệ quy & $\Theta(N^2)$ & $\Theta(N^2)$\\
  Product hai chiều & $\Theta(N^{3/2})$ & $\Theta(N^2)$\\
  \hline
 \end{tabular}
\end{table}

Trường hợp $N=1$ được xử lý riêng: chỉ có một biến ô, không có
clause AMO, và chương trình giữ hai clause ALO đơn vị trùng nhau của
hàng và cột. Số lượng là một biến, hai clauses. Đối với $N\geq2$,
các công thức trên khớp với cách duyệt nhóm của chương trình. So sánh
thực nghiệm cần dùng số đếm đã sinh thực tế, đồng thời dùng công thức
để phát hiện sai khác trong cài đặt.

\section{Các mô hình đối chứng CP và MIP}
\label{sec:theory-baselines}

\subsection{Mô hình CP và CP-SAT}

Lập trình ràng buộc (\emph{constraint programming}, CP) mô tả trực tiếp
miền giá trị và các quan hệ cần thỏa. Chương trình dùng biến nguyên
$q_r\in\{0,\ldots,N-1\}$ và ba ràng buộc
\begin{align}
 &\operatorname{AllDifferent}(q_0,\ldots,q_{N-1}),\label{eq:th-cp-column}\\
 &\operatorname{AllDifferent}(q_0+0,\ldots,q_{N-1}+N-1),\\
 &\operatorname{AllDifferent}(q_0-0,\ldots,q_{N-1}-(N-1)).
 \label{eq:th-cp-diagonal}
\end{align}
\emph{AllDifferent} yêu cầu các biểu thức trong danh sách nhận các
giá trị đôi một khác nhau. Hai hậu chung đường chéo khi và chỉ khi
$q_r+r=q_t+t$ hoặc $q_r-r=q_t-t$, nên ba ràng buộc này tương đương
\eqref{eq:th-columns}--\eqref{eq:th-diagonal}. Mô hình có $N$ biến
nguyên ở tầng biểu diễn của chương trình; số biến nội bộ sau chuyển
đổi của solver là đại lượng khác.

Cùng mô hình được xây dựng cho OR-Tools CP-SAT \cite{ortools} và
IBM CP Optimizer thông qua DOcplex CP \cite{ibm}. Đây là hai bộ giải
khác nhau; việc chạy CP-SAT không được ghi thành kết quả CPLEX CP.
Không bổ sung mục tiêu tối ưu, vì nhiệm vụ chỉ cần một nghiệm hợp lệ.
Các hàm đối chứng đặt một luồng tìm kiếm và hạt giống ngẫu nhiên
(\emph{random seed}) bằng 42; các cơ chế tiền xử lý và xử lý đối xứng
nội bộ khác giữ theo mặc định của solver.

\subsection{Mô hình nguyên nhị phân}

Quy hoạch nguyên tuyến tính (\emph{integer linear programming}, ILP)
được dùng dưới dạng mô hình nhị phân và giải bằng CPLEX MIP
\cite{ibm} hoặc Gurobi \cite{gurobi}. MIP
(\emph{mixed-integer programming}) là tên họ bộ giải; mô hình cụ thể
ở đây chỉ cần các biến nhị phân $z_{r,c}\in\{0,1\}$:
\begin{align}
 \min\quad &0,\\
 \text{s.t.}\quad
 &\sum_{c=0}^{N-1}z_{r,c}=1 &&0\leq r<N,\label{eq:th-mip-row}\\
 &\sum_{r=0}^{N-1}z_{r,c}=1 &&0\leq c<N,\label{eq:th-mip-column}\\
 &\sum_{r-c=d}z_{r,c}\leq1 &&|D^-_d|\geq2,\\
 &\sum_{r+c=s}z_{r,c}\leq1 &&|D^+_s|\geq2.
 \label{eq:th-mip-diag}
\end{align}
Trong các tổng đường chéo, chỉ lấy các ô thuộc bàn cờ. Mô hình có
$N^2$ biến nhị phân và, với $N\geq2$, $6N-6$ ràng buộc tuyến tính
không kể khai báo miền. Các phương trình hàng/cột bảo đảm đúng $N$
hậu; bất đẳng thức đường chéo loại các cặp tấn công. Vì vậy mô hình
biểu diễn cùng tập bàn cờ hợp lệ như \eqref{eq:th-board}. Mục tiêu
hằng 0 cho phép dùng giao diện tối ưu nhưng không làm thay đổi nhiệm
vụ khả thi (\emph{feasibility problem}).

\subsection{Điều kiện so sánh và diễn giải trạng thái}

Các bộ giải đều dừng ở một lời giải hợp lệ hoặc một kết luận không
khả thi, với giới hạn thời gian được cấu hình trong quy trình thực
nghiệm. Một nghiệm được trả về từ CP hay MIP phải qua kiểm tra hình
học độc lập giống nghiệm SAT. Đối với MIP, chương trình đọc các giá
trị lớn hơn $0.5$ là ô có hậu rồi kiểm tra lại toàn bộ bàn cờ; kết
quả không vượt qua kiểm tra không được tính là nghiệm hợp lệ.

Trạng thái bản địa của công cụ được giữ lại và ánh xạ sang các nhóm
chung: SAT khi có nghiệm đã kiểm tra, UNSAT khi công cụ kết luận không
khả thi, TIMEOUT khi chưa có kết luận trong giới hạn, UNAVAILABLE khi
thiếu công cụ, LICENSE\_LIMIT khi bị giới hạn giấy phép, và ERROR cho
lỗi khác. Đặc biệt, thư viện mô hình DOcplex được cài đặt không có
nghĩa là bộ thực thi CP Optimizer đã sẵn sàng. Các lần chạy thiếu công
cụ hoặc bị chặn bởi giấy phép phải được công bố riêng, không xem là
UNSAT hay giả lập bằng một bộ giải khác.

Số biến và ràng buộc của CP, MIP và CNF được xác định ở những tầng
biểu diễn khác nhau. So sánh hiệu năng vì thế dựa trên trạng thái đã
xác nhận, chi phí xây dựng, chi phí giải và tổng thời gian theo phạm
vi đo đã công bố. Phần thực nghiệm trình bày các chỉ số này cùng
giới hạn tài nguyên và cấu hình công cụ để xác định phạm vi áp dụng
của kết luận.

% Gọi \TheorySourceAppendix sau \appendix trong tệp báo cáo chính.
% Định nghĩa ở đây chỉ giữ nội dung nguồn của fragment; không đổi bộ đếm
% section trước khi main.tex chèn thực nghiệm và kết luận.
\newcommand{\TheorySourceAppendix}{%
\section{Phạm vi và đối chiếu tài liệu nguồn}
\label{app:theory-sources}

\subsection{Phạm vi tài liệu đã đọc}

Mục này tóm tắt cách sử dụng các tài liệu đầu vào. Phạm vi đã đọc là
toàn bộ phần văn bản được trích xuất theo trang PDF của năm tài liệu
dưới đây; số trang tính từ trang đầu tệp, bao gồm bìa. Hình được đối
chiếu qua nhãn và chú thích trong bản trích, không số hóa đường cong
trong hình thành dữ liệu thực nghiệm.
\begin{enumerate}
 \item \emph{Assignment 1 -- requirements.pdf}, trang 1/1:
       xác định đề tài, yêu cầu cài đặt, đối chứng và công bố dữ liệu.
 \item \emph{An optimized SAT approach solving N-Queens problem.pdf},
       trang 1--41/41: tham khảo mô hình bàn cờ, năm mã hóa AMO và
       cách tổ chức thực nghiệm \cite{bach2024}.
 \item \emph{Mau xay dung bao cao BTL so 1 và BTL so 2.pdf},
       trang 1--32/32: tham khảo bố cục bài báo, trình bày phương pháp,
       thực nghiệm, kết luận, phụ lục và Data Availability.
 \item \emph{Module3 AMO ALO Encoding}, trang 1--17/17:
       đối chiếu định nghĩa và cách xây dựng các mã hóa AMO/ALO
       \cite{module3}.
 \item \emph{Module2 SAT Encoding}, trang 1--27/27:
       tham khảo CNF, biến phụ và tính đồng khả thỏa.
\end{enumerate}
Bảng thực nghiệm trong các tài liệu nguồn được xem là kết quả công
bố của tài liệu tương ứng. Dữ liệu trong phần thực nghiệm của báo
cáo được lấy từ các lần chạy của chương trình đi kèm.

\subsection{Đối chiếu công thức}

Các công thức được viết lại nhất quán với cài đặt và chứng minh trong
Mục~\ref{sec:theory-method}. Đối với \cite{bach2024}, ví dụ ALO ở
trang PDF 15--16 được chuẩn hóa thành một clause cho từng hàng hoặc
cột; các đường chéo chỉ có AMO. Clause cuối của sequential ở trang
PDF 32--33 dùng dấu phủ định của biến ô, còn các liên kết product ở
trang PDF 26, 34 được thống nhất chỉ số biến chọn. Việc đối chiếu này
nhằm làm rõ phần trình bày công thức trong nguồn; mã thực nghiệm của
tác giả nguồn không thuộc phần được kiểm tra.
}
